\scanpage{52}
% archive sheet number: 51

\mspage{16}
Dans Hot, par contre on a ceci :

Tout carré cocartésien est cartésien. Pour qu'un carré cartésien

\diagram{\begin{tikzcd}[ampersand replacement=\&]
\& \xi_0\arrow[dl,no head]\arrow[dr,no head] \&\\
\xi_a\arrow[dr,no head] \& \& \xi_b\arrow[dl,no head]\\
\& \xi_1 \&
\end{tikzcd}}

soit aussi cocartésien, il f. et s. que l'on ait

\[
\pi_0(\xi_1)=\mathrm{Im}\,\pi_0(\xi_a)\cup\mathrm{Im}\,\pi_0(\xi_b)
\]

% slanted
\margin{complètement faux. Si p.ex. $\xi_1=e$, le diagramme ci-contre : le fait homotopiquement cartésien et cocartésien pourvu que $\xi_a\simeq e$ ou $\xi_b\simeq e$}

Der 6 ou Der 6' au choix donnent W7. Quant à W8, il semble tributaire d'un principe général \inserted{(``homom. des cinq'')}

``\uline{Der 8}'' Soit un hom. d'un carré \deleted{\unsure{bi}}cartésien dans un autre, dans $A^{\prime\prime}(\Delta^1\times\Delta^1)$. Considérons

\diagram{\begin{tikzcd}[ampersand replacement=\&,column sep=small,row sep=small]
\& \xi_a\arrow[dr,no head]\arrow[ddd,"f_a"'{pos=.22}] \& \&\\
\xi_0\arrow[ur]\arrow[dr]\arrow[ddd,"f_0"'] \& \& \xi_1\arrow[ddd,"f_1"] \&\\
\& \xi_b\arrow[ur]\arrow[ddd,"f_b"] \& \&\\
\& \eta_a\arrow[dr] \& \&\\
\eta_0\arrow[ur]\arrow[dr] \& \& \eta_1 \&\\
\& \eta_b\arrow[ur] \& \&
\end{tikzcd}}

les conditions

a) $f_0$ est iso

b) $f_1$ est iso

c) $f_a$ \uline{et} $f_b$ est iso.

Alors \inserted{si les carrés sont ``isom.'' cartésiens,} \deleted{\unsure{chacune des deux}} $b+c\Rightarrow a$. S'ils sont bien cocartésiens, $a+c\Rightarrow b$. \deleted{\unsure{Si les carrés sont cartésiens}} S'ils sont bicartésiens, on a \deleted{les formules} \unsure{(il est)} de plus $a+b\Rightarrow c$ (?).

\margin{\inserted{vérifier si c'est raisonnable, pas évident}}

% slanted
\margin{Il suffit d'exiger que $a)+b)\Rightarrow c$ i.e. si $f_0,f_1$ iso, alors $f_a$ et $f_b$ iso. \unsure{Ceci en|Tout ceci en} évidence.}

\uuline{NB} \deleted{La \unsure{form.} $a+c\Rightarrow b$, \unsure{est|étant} cocartésien donc $b+c\Rightarrow a$ s'il suffit cartésien et c'est évident. Le cas non évident est le cas $a+b\Rightarrow c$}

% slanted
\margin{version \unsure{renforcée} de W7 page 4.}

\mspage{17}
\uline{W7bis} Soit

\diagram{\begin{tikzcd}[ampersand replacement=\&,row sep=small]
\& X_a\\
X_0\arrow[ur]\arrow[dr] \&\\
\& X_b
\end{tikzcd}}

diagramme dans Cat, soit

$\underline\Phi={}$\diagram{\begin{tikzcd}[ampersand replacement=\&,row sep=small]
\& a\arrow[dl]\\
0 \&\\
\& b\arrow[ul]
\end{tikzcd}}

et $\mathcal X$ la catégorie fibrée sur $\underline\Phi$ définie par le diagramme. On a un diagramme commutatif dans Hot$_W$ :

\diagram{\begin{tikzcd}[ampersand replacement=\&]
\& X_0\arrow[dl,"i_a"']\arrow[dr,"i_b"] \&\\
X_a\arrow[dr,"i_{b'}"'] \& \& X_b\arrow[dl,"i_{a'}"]\\
\& \mathcal X \&
\end{tikzcd}}

% slanted
\margin{\uline{NB} Autre présentation. On \uline{part} de $\mathcal X$, fibrée sur}

\diagram{\begin{tikzcd}[ampersand replacement=\&,row sep=tiny,column sep=tiny]
\& a\arrow[dl]\\
0 \&\\
\& b\arrow[ul]
\end{tikzcd}}

(NB Dans Cat, \deleted{je présume que} les deux flèches composées sont homotopes, par un chemin de longueur 2 \deleted{un} elles sont homotopes tous deux, par une homotopie de long. 1, à ``l'inclusion'' $X_0\lhook\joinrel\longrightarrow\mathcal X$)

Alors \deleted{\unsure{iso|tous}}

\[
\begin{aligned}
i_a\in W&\Longleftrightarrow i_{a'}\in W\\
i_b\in W&\Longleftrightarrow i_{b'}\in W
\end{aligned}
\]

\uline{W8bis} Soit un diagramme comm. dans Cat

\diagram{\begin{tikzcd}[ampersand replacement=\&,column sep=small,row sep=small]
\& X_a\arrow[dr,dashed]\arrow[ddd,"f_a"'{pos=.22}] \&\\
X_0\arrow[ur]\arrow[dr]\arrow[ddd,"f_0"'] \& \& \mathcal X\arrow[ddd,"f"]\\
\& X_b\arrow[ur,dashed]\arrow[ddd,"f_b"] \&\\
\& X'_a\arrow[dr,dashed] \&\\
X'_0\arrow[ur]\arrow[dr] \& \& \mathcal X'\\
\& X'_b\arrow[ur,dashed] \&
\end{tikzcd}}

\margin{(traits pleins)}

% slanted
\margin{équivalent : l'énoncé \unsure{direct}, avec $\mathcal X,\mathcal X'$ cofibrées sur $\underline\Psi$ et un \uline{cocartésien} $\mathcal X\to\mathcal X'$.}

% slanted
\margin{à examiner, ici et dans W7bis, si l'hyp. ``$\unsure{f'|f}$ cartésien'' (resp. cocartésien) est nécessaire ou superflue ?}

On le complète par un diagramme commutatif dans Hot$_W$ (``par construction'' de $\mathcal X,\mathcal X'$\ \deleted{dans}, où pour

\diagram{$X=\displaystyle\int\left(\begin{tikzcd}[ampersand replacement=\&,row sep=small,column sep=small]
\& X_a\\ X_0\arrow[ur,no head]\arrow[dr] \&\\ \& X_b
\end{tikzcd}\right)$}

). Considérons les cond.

\[
\left\{\begin{aligned}
\text{a)}&\quad f_{\unsure{*|1}}\in W\\
\text{b)}&\quad f_0\in W\\
\text{c)}&\quad f_a\ \text{et}\ f_b\in W
\end{aligned}\right.
\]

Alors deux \deleted{\unsure{deux|des}} ces conditions impliquent la troisième.
